
%%%%%%%%%%%%%%%%%%%%%%%%%%%%%%%%%%%%%%%%%%%%%%%%%%%%%%%%%%%%%%%%%%%%%%%%%%%%%%%%%%%%%%%%%%%%%%%%%%% 
% KAPITOLA 7
%%%%%%%%%%%%%%%%%%%%%%%%%%%%%%%%%%%%%%%%%%%%%%%%%%%%%%%%%%%%%%%%%%%%%%%%%%%%%%%%%%%%%%%%%%%%%%%%%%
\chapter{Metodika experimentov}
\label{ch:kap7}

Kapitola~\ref{ch:kap6} opísala implementáciu nástroja a táto kapitola definuje spôsob, akým sa vyhodnocuje kvalita generovaných receptov. Cieľom tejto kapitoly nie je zhodnotenie správnosti kódu. Kapitola najskôr formuluje ciele experimentov(Sekcia ~\ref{sub:ciele}) a definuje testovacie databázy (Sekcia~\ref{sub:testovacie_databazy}) a vrstvy vyhodnotenia ako aj celkové spôsoby testovania (Sekcia~\ref{sub:sposob_vyhodnotenia_uspesnosti}). Samotné výsledky sú prezentované v Kapitole~\ref{ch:kap8}.

\section{Ciele}
\label{sub:ciele}
Účel testovania je zodpovedať niekoľko otázok, ktoré pokrývajú konfiguráciu LLM, vlastnosti vstupných schém a spoľahlivosť vyhodnocovacej metodiky.

\begin{itemize}
    \item \textbf{RQ1} Aký je vplyv teploty vzorkovania na úspešnosť pipeline (L1-L3)?
    \item \textbf{RQ2} Aký je vplyv veľkosti modelu v rámci jednej rodiny na úspešnosť?
    \item \textbf{RQ3} Líši sa úspešnosť medzi modelovými rodinami porovnateľnej veľkosti?
    \item \textbf{RQ4} Ako sa mení úspešnosť s rastúcou zložitosťou databázovej schémy?
    \item \textbf{RQ5} Pri akej veľkosti schémy dochádza k poklesu kvality a úspešnosti ?
    \item \textbf{RQ6} Je možné pri vyhodnocovaní spoľahlivo odlíšiť zlyhania spôsobené navrhnutou pipeline od zlyhaní generátora Dataster?
\end{itemize}

Celkovo je na testovanie pripravených 19 základných testovacích databáz (testdb-testdb20), 3 testovacie databázy na reprodukciu chýb a 9 extrémnych testovacích databáz (extreme\_testdb-extreme\_testdb9). Každá z nich slúži na testovanie pipeline na rôznych typoch  databázových schém. Databáza testdb13 bola zo sady vylúčená ako duplikát testdb12.

Databázy sú usporiadané podľa rastúcej zložitosti. Samotná zložitosť je definovaná pomocou týchto kritérií:
\begin{itemize}
    \item \textbf{počet tabuliek a stĺpcov:} s rastúcim počtom sa zvyšuje objem dát, ktoré musí LLM spracovať
    
    \item \textbf{počet cudzích kľúčov:} miera prepojenia medzi tabuľkami 
    
    \item \textbf{zložitosť primárnych kľúčov:} prítomnosť zložených primárnych kľúčov, ktoré komplikujú jednoznačnú identifikáciu a prepojenie záznamov

\end{itemize}


\section{Testovacie databázy}
\label{sub:testovacie_databazy}

\subsection{Základné testovacie databázy}

Účelom týchto testovacích dát je skontrolovať celkovú funkčnosť programu.

\begin{table}[H]
    \centering
    \small
    \begin{tabular}{>{\centering\arraybackslash}p{0.1\linewidth}>{\centering\arraybackslash}p{0.1\linewidth}>{\centering\arraybackslash}p{0.1\linewidth}>{\centering\arraybackslash}p{0.08\linewidth}>{\centering\arraybackslash}p{0.08\linewidth}>{\centering\arraybackslash}p{0.4\linewidth}}
    \toprule
    \textbf{Názov db} & \textbf{počet tabuliek} & \textbf{počet stĺpcov} & \textbf{počet FK} & \textbf{zložený PK} & \textbf{Cieľ testovania} \\
    \midrule
         testdb   & 1 & 3  & 0 & nie & zákl. funkcionalita (INT+VARCHAR)  \\
         testdb2  & 1 & 3  & 0 & nie & nízky počet unikátnych hodnôt  \\
         testdb3  & 1 & 7  & 0 & nie & desatinné čísla, prázdne dátumy  \\
         testdb4  & 1 & 9  & 0 & nie & kontrola významu stĺpcov  \\
         testdb5  & 1 & 10 & 0 & nie & NULL hodnoty, poradie dátumov  \\
         testdb6  & 1 & 10 & 0 & nie & kontrola významu reťazcov  \\
         testdb7  & 1 & 4  & 0 & nie & nízky počet unikátnych hodnôt, diakritika \\
         testdb8  & 1 & 6  & 0 & nie & INT rozsah  \\
         testdb9  & 1 & 5  & 0 & nie & regex (vzorové reťazce INV-YYYY-NN)  \\
         testdb10 & 2 & 15 & 1 & nie & FK vzťah  \\
         testdb11 & 2 & 16 & 1 & nie & CamelCase/Snake\_Case, NULL hodnoty  \\
         testdb12 & 2 & 18 & 1 & nie & FK Courses $\rightarrow$ Authors\_Tutors\\
         testdb14 & 2 & 13 & 1 & nie & non-numeric PK  \\
         testdb15 & 3 & 26 & 2 & nie & 1:N  \\
         testdb16 & 3 & 8  & 2 & áno & zložený PK  \\
         testdb17 & 2 & 11 & 0 & nie & testovanie constraints  \\
         testdb18 & 5 & 15 & 5 & nie & hĺbka FK reťazca (4 úrovne) \\
         testdb19 & 3 & 7  & 2 & áno & M:N  \\
         testdb20 & 1 & 4  & 0 & nie & CamelCase, NULL hodnoty   \\
         \bottomrule
         
    \end{tabular}
    \caption{Základné testovacie databázy}
    \label{tab:zaklad_databazy}
\end{table}



\subsection{Testovacie databázy na reprodukciu chýb}
Tieto testovacie databázy slúžia na reprodukciu obmedzenia generátora Dataster pri prevode hodnôt zo SMT riešiteľa do SQL pre cieľovú DBMS (MySQL). Cieľom nie je hodnotiť Dataster, ale overiť, že vyhodnocovacia metodika dokáže tieto prípady správne odlíšiť od zlyhaní pipeline.

\begin{table}[H]
    \centering
    \small
    \begin{tabular}{>{\centering\arraybackslash}p{1.5cm}>{\centering\arraybackslash}p{2cm}>{\centering\arraybackslash}p{1.5cm}>{\centering\arraybackslash}p{1.5cm}>{\centering\arraybackslash}p{1.5cm}>{\centering\arraybackslash}p{5cm}}
    \toprule
    \textbf{Názov db} & \textbf{počet tabuliek} & \textbf{počet stĺpcov} & \textbf{počet FK} & \textbf{zložený PK} & \textbf{Cieľ testovania} \\
    \midrule
         testdb21 & 1 & 4 & 0 & nie & DECIMAL/float  \\
         testdb22 & 1 & 3 & 0 & nie & SMT uniqueness  \\
         testdb23 & 1 & 4 & 0 & nie & DATE/DATETIME/TIME  \\
         \bottomrule
         
    \end{tabular}
    \caption{Testovacie databázy na reprodukciu chýb}
    \label{tab:reprodukcia_chyb_databazy}
\end{table}


\subsection{Extrémne testovacie databázy}

Zatiaľ čo predošlé testy kontrolujú či program funguje bez problémov, tieto testovacie sady boli vytvorené pre zaťaženie programu na hranici jeho limitov. Ich cieľom je zistiť ako pri vysokom počte tabuliek, šírke tabuliek a hĺbke/tvare FK grafu klesá úspešnosť alebo kvalita výsledkov.

\begin{table}[H]
    \centering
    \small
    \begin{tabular}{>{\raggedright\arraybackslash}p{3cm}>{\centering\arraybackslash}p{1cm}>{\centering\arraybackslash}p{1cm}>{\centering\arraybackslash}p{1cm}>{\centering\arraybackslash}p{1.2cm}>{\raggedright\arraybackslash}p{6cm}}
    \toprule
    \textbf{Názov db} & \textbf{počet tabuliek} & \textbf{počet stĺpcov} & \textbf{počet FK} & \textbf{zložený PK} & \textbf{Cieľ testovania} \\
    \midrule
         extreme\_testdb & 30 & 150 & 29 & nie & škálovanie počtu tabuliek  \\
         extreme\_testdb2 & 60 & 300 & 59 & nie & škálovanie počtu tabuliek  \\
         extreme\_testdb3 & 100 & 500 & 99 & nie & škálovanie počtu tabuliek  \\
         extreme\_testdb4 & 1 & 50 & 0 & nie & ako šírka tabuľky zaplní kontextové okno  \\
         extreme\_testdb5 & 1 & 100 & 0 & nie & ako šírka tabuľky zaplní kontextové okno  \\
         extreme\_testdb6 & 1 & 200 & 0 & nie & ako šírka tabuľky zaplní kontextové okno  \\
         extreme\_testdb7 & 10 & 19 & 9 & nie & hĺbka FK reťazca  \\
         extreme\_testdb8 & 20 & 39 & 19 & nie & hĺbka FK reťazca   \\
         extreme\_testdb9 & 21 & 41 & 20 & nie & fan-out   \\
         \bottomrule
         
    \end{tabular}
    \caption{Extrémne testovacie databázy}
    \label{tab:extremne_databazy}
\end{table}


\section{Spôsob vyhodnotenia úspešnosti}
\label{sub:sposob_vyhodnotenia_uspesnosti}
Proces vyhodnotenia úspešnosti vygenerovaných receptov prebieha na troch úrovniach, kde každá z nich definuje špecifické kritériá pre úspech alebo zlyhanie

\begin{figure}[H]
    \centering
    \includegraphics[width=0.5\linewidth]{obrazky-figures/flowtestovania.pdf}
    \caption{Trojúrovňoví proces vyhodnotenia recepu}
    \label{fig:flowtestovania}
\end{figure}

\subsection{Úroveň L1: overenie správnosti YAML}
\label{sub:uroven_l1_overenie_spravnost}
Táto úroveň sa zameriava na úspešnosť pipeline a kontrolu štruktúry vygenerovaného receptu.  Táto vrstva overuje formát YAML,  výskyt povinných sekcií a dodržiavanie pravidiel validačnej vrstvy kódu. Celková úspešnosť sa počíta ako podiel úspešných receptov k celkovému počtu vygenerovaných receptov:

\begin{equation}
        L1\_result = \frac{L1\_success}{L1\_overall}    
\end{equation}

\begin{description}
    \item[L1\_success] počet receptov, ktoré prešli validačnou vrstvou
    \item[L1\_overall] celkový počet vygenerovaných receptov
\end{description}

\subsection{Úroveň L2: akceptácia nástrojom Dataster}
Ďalšia úroveň kontroluje, či je nástroj Dataster schopný úspešne spracovať recept a vygenerovať testovaciu sadu dát. Overuje akceptáciu Datasterom, a preto ak je objavená chyba, je to okamžite vyhodnotené ako zlyhanie. L2 sa vyhodnocuje nad všetkými behmi, v ktorých vznikol recept, výsledok L1 sa hodnotí samostatne a recepty sa pre L2 sa nefiltrujú.

\begin{equation}
        L2\_result = \frac{L2\_success}{L2\_overall}
\end{equation}

\begin{description}
    \item[L2\_success] počet receptov, ktoré boli akceptované nástrojom Dataster
    \item[L2\_overall] celkový počet behov s vygenerovaným receptom
\end{description}

\subsection{Úroveň L3: exekúcia SQL dotazov a kontrola databázy}

Posledná vrstva kontroluje použiteľnosť vygenerovaných SQL dotazov v databáze. Na rozdiel od ostatných vrstiev sa tu jednotka meranie mení z receptov na jednotlivé SQL príkazy. Príkazy sú spustené v naklonovanej databáze bez dát (obsahujú iba pôvodnú schému klonovanej databázy). Na tejto vrstve sa kontroluje, či tvary vygenerované programom Dataster sú akceptované vybranou DBMS (MySQL). Cieľom je nestranne zmerať úspešnosť vykonanie SQL výstupu, pričom zlyhania môžu byť spôsobené buď chybou v kóde, alebo zdokumentovaným obmedzením samotného nástroja. Každý neakceptovaný príkaz je zaznamenaný spolu s klasifikáciu zlyhania. Pre kategóriu dataster\_sql\_escaping je vyhodnotenie dodatočne vykonané v odlišnom režime.

\begin{equation}
       L3\_raw = \frac{L3\_success}{L3\_overall}
\end{equation}

\begin{equation}
        L3\_adjusted = \frac{L3\_success}{L3\_overall - L3\_dataster\_fail}
\end{equation}

\begin{description}
    \item[L3\_success] počet úspešne vykonaných SQL dotazov
    \item[L3\_overall] celkový počet vygenerovaných SQL dotazov
    \item[L3\_dataster\_fail] počet INSERT príkazov, ktoré boli klasifikované ako obmedzenia nástroja Dataster
\end{description}


\section{Klasifikácia zlyhaní}

Pri vyhodnocovaní úspešnosti je nutné jasne určiť, ktorá časť systému zlyhanie spôsobila. Táto tabuľka ukazuje presné pravidlá, podľa ktorých sa každé zlyhanie triedi. Vďaka tomu vieme presne vypočítať metriku L3\_Adjusted.

\begin{table}[H]
    \centering
    \small
    \begin{tabular}{>{\raggedright\arraybackslash}p{0.25\linewidth}>{\centering\arraybackslash}p{0.12\linewidth}>{\raggedright\arraybackslash}p{0.4\linewidth}>{\centering\arraybackslash}p{0.15\linewidth}}
    \toprule
    \textbf{ID zlyhania} & \textbf{úroveň} & \textbf{podmienka} & \textbf{komu patrí zlyhanie} \\
    \midrule
        pipeline\_invalid\_yaml & L1 & výstup nie je platný YAML alebo chýba povinná sekcia alebo recept nevznikol (timeout, cyklické FK, prázdna taabuľka) & pipeline \\
         
        dataster\_smt\_count & L2 & generátor nedokáže vytvoriť požadovaný počet unikátnych záznamov (GenerateCountError) & Dataster \\
         
        zero\_rows & L2 & recept prijatý, ale vygenerovaných 0 riadkov & pipeline \\

         dataster_recursion & L2 & RecursionError pri extrémne tabuľkách & Dataster\\
         
         iná výnimka & L2 & Dataster odmietol recept z iného dôvodu & pipeline \\
         
        dataster\_sql\_escaping & L3 &  syntaktická chyba SQL (errno 1064) spôsobená neescapovanými znakmi v hodnotách & Dataster \\
         
        dataster\_finterval\_fraction & L3 & odmietnutá hodnota v tvare zlomku pre stĺpec DECIMAL (errno 1366) & Dataster \\
         
        pipeline\_type\_mismatch & L3 &  odmietnutá hodnota pre nesúlad typov, bez zlomku (errno 1366) & pipeline \\
        
        dataster\_cascade\_fk & L3 & porušenie FK (errno 1452), rodičovský riadok zlyhal na dataster\_* chybe  & Dataster \\
        
        fk\_violation & L3 & porušenie FK (errno 1452), táto chyba nie je len reakciou na zlyhanie rodiča, ale má vlastnú príčinu & pipeline\\
        
        duplicate\_pk & L3 & porušenie unikátneho kľúča (errno 1062) &  pipeline\\
        
        not\_null\_violation & L3 & NULL v NOT NULL stĺpci (errno 1048) & pipeline \\
        
        data\_too\_long & L3 &  hodnota presahuje dĺžku stĺpca (errno 1406) & pipeline\\
        
        other & L3 & iná chyba MySQL & pipeline\\ 
        \bottomrule
        
    \end{tabular}
    \caption{Klasifikácia zlyhaní}
    \label{tab:klasifikacia_zlyhani}
\end{table}

Zlyhanie, ktoré nie je možné jednoznačne priradiť zdokumentovanému obmedzeniu nástroja Dataster, sú ponechané ako zlyhanie pipeline. Metrika L3\_Adjusted teda vylučuje iba kategórie s prefixom dataster\_

\section{Testovacie scenáre a ich vykonanie}
\label{sub:testovacie_scenare}

Testovanie prebieha v niekoľkých fázach. Hlavne sa program testuje na referenčnom DBMS MySQL, ale keďže je implementované SQLAlchemy tak pre demonštráciu je vykonaný program s DBMS PostgreSQL s podmnožinou základných testovacích databáz. Jadrom testovanie je model qwen2.5 s rôznymi teplotami a potom nasleduje testovanie (Rozšírenie) s rôznymi rodinami pre účel porovnania výkonu modelov. Pri testovaní sa používajú testovacie databázy spomenuté v Sekcii~\ref{sub:ciele}.

\begin{figure}[h]
    \centering
    \includegraphics[width=1\linewidth]{obrazky-figures/rozdelenietestovania.pdf}
    \caption{Tri fázy testovania}
    \label{fig:rozdelenie_testovania}
\end{figure}

Každá konfigurácia beží v troch kolách, výstupy sa ukladajú oddelene pre každé kolo a analyzujú oddelene

\subsection{Jadro}
V jadre testovania sa zameriavame výhradne na model qwen2.5 a spúšťame program naprieč rôznymi teplotami pre rôzne veľkosti parametrov modelu, ako 3b, 7b a mix predom spomenutých dvoch modelov.

\begin{table}[H]
    \centering
    \small
    \begin{tabular}{lccc}
    \toprule
    
    \textbf{Model} & \textbf{teplota} & \textbf{počet kôl} & \textbf{počet DB} \\
    \midrule
        qwen2.5:3b/qwen2.5:7b & {0.0, 0.2, 0.3, 0.5, 0.8, 1.2} & 3 & 31 \\
        qwen2.5:3b & {0.0, 0.2, 0.3, 0.5, 0.8, 1.2} & 3 & 31 \\
        qwen2.5:7b & {0.0, 0.2, 0.3, 0.5, 0.8, 1.2} & 3 & 31 \\ \bottomrule
        
               
    \end{tabular}
    \caption{jadro testovania}
    \label{tab:jadro_testovania}
\end{table}

\subsection{Rozšírenie}

Ďalej pokračujeme na rozšírenie, kde sa zameriava na samotné rodiny modelov. Cieľom je porovnať výsledky medzi rôznymi rodinami ako aj mix týchto modelov. V tomto prípade sa nezameriavame na rôzne teploty, ale použivame referenčné nastavenie teploty (0.0).

\begin{table}[H]
    \centering
    \small
    \begin{tabular}{lc>{\centering\arraybackslash}p{0.1\linewidth}l}
    \toprule
    
    \textbf{Model} & \textbf{teplota} & \textbf{počet kôl} & \textbf{počet DB} \\
    \midrule
        llama3.1:8b & referenčná teplota 0.0 & 3 & 19 (\ref{tab:zaklad_databazy})\\
        mistral:7b & referenčná teplota 0.0 & 3 & 19 (\ref{tab:zaklad_databazy}) \\ \bottomrule         
    \end{tabular}
    \caption{rozšírenie testovania}
    \label{tab:rozsirenie_testovania}
\end{table}

\subsection{Demonštrácia}
Ako posledné je demonštrácia použitia iného DBMS (PostgreSQL), teda, či implementácia funguje aj pre inú DBMS ako referenčnú.

\begin{table}[H]
    \centering
    \small
    \begin{tabular}{cccc}
    \toprule
    
    \textbf{Model} & \textbf{teplota} & \textbf{počet kôl} & \textbf{DB} \\
    \midrule
        qwen2.5:3b/qwen2.5:7b & 0.0 & 3 & testdb - testdb10\\ \bottomrule          
    \end{tabular}
    \caption{demonštrácia testovania}
    \label{tab:demonstracia_testovania}
\end{table}

\section{Hranice vyhodnotenia}

\begin{itemize}
    \item \textbf{Verzia softvéru:} prezentované výsledky sú naviazané na konkrétne verzie DBMS a verziu programu Dataster. Pri využití iných verzii sa správanie programu môže líšiť. 
    \item \textbf{Počet testovacích kôl:} použitie troch kôl je pre ukážku rozptylu výsledkov, nie je to však postačujúce na plnohodnotné testovanie.
    \item \textbf{Hardvérové limity:} výber testovacích modelov je obmedzený dostupným hardvérom a výsledky neplatia pre väčšie a výkonnejšie modely.
    \item \textbf{Charakter dát:} testovacie databázy slúžia na pokrytie rôznych testovacích prípadov a teda nereprezentujú skutočné dáta z reality. Nameranú úspešnosť preto nie je možné brať ako istotu pre reálne nasadenie.
\end{itemize}

